\cotefile{64}
\scanpage{104}\mspage{2}
On trouve
\[
 (\Theta_\delta f)^\vee=\Theta_{\check{\delta}}\check f
 =\Theta_{-\delta}\check f=-\Theta_\delta\check f \tag*{(15)}
\]
i.e.
\[
 v\circ\Theta_\delta=-\Theta_\delta\circ v\,.
\]
Posons
\[
 \left\{\begin{aligned}
 F_n^+&=\{f\in F_n\mid\check f=f\}\\
 F_n^-&=\{f\in F_n\mid\check f=-f\}
 \end{aligned}\right. \tag*{(16)}
\]
on a donc
\[
 \left\{\begin{aligned}
 &\deleted{\unsure{\{}}\,F_n=F_n^+\oplus F_n^-\\
 &\left\{\begin{aligned}
 &\Theta_\delta F_n^+\subset F_n^-,\quad \Theta_\delta F_n^-\subset F_n^+\\
 &t_E\otimes F_n^+\xrightarrow{\delta_n}F_n^-,\quad t_E\otimes F_n^-\xrightarrow{\delta_n}F_n^+
 \end{aligned}\right.\\
 &\left\{\begin{aligned}
 F_n^+F_m^+&\subset F_{n+m}^+\\
 F_n^{\unsure{+|-}}F_m^-&\subset F_{n+m}^+
 \end{aligned}\right.\qquad F_n^+F_m^-\subset F_{n+m}^-
 \end{aligned}\right. \tag*{(17)}
\]
On va construire \inserted{l'alg. graduée} $F_*$, avec la structure $\delta_*:t_E\otimes F_*\mathrel{\unsure{-|\to}}F_*$, \inserted{i.e $t_E\to\operatorname{Der}^{+1}_K(F_*,F_*)$,} et la structure $v$ (automorphisme d'ordre 2).

On a
\[
 F_2^+=F_2\xleftarrow{\sim}\Gamma(\Sigma_{J},\{t\})
 \qquad\text{où}\quad
 \begin{aligned}
 J&={}_2E^*\\
 \Sigma_J&\simeq E/(\pm\operatorname{id}_E)\\
 &t\ \text{image de }0_E\\
 &t\in\Sigma_J^*=\Sigma_J\smallsetminus J
 \end{aligned} \tag*{(18)}
\]
plus généralement on a
% archive sheet number: 103
\scanpage{105}
\[
 F_{2n}^+\xrightarrow{\sim}F_{2n+1}^+=\Gamma(\Sigma_{J},n\{t\}) \tag*{(19)}
\]
On a d'autre part $F_n^-=0$ si $n\leq 2$, et
\[
 F_3^-=\delta_2(t_E\otimes F_2)\overset{\unsure{\varepsilon}}{\simeq}t_E^{\otimes3} \tag*{(20)}
\]
supplémentaire de rang 1 de $F_2=F_3^+$ dans $F_3$.

On a alors, pour \deleted{$n\geq3$ \unsure{de|et|tout}} $n\geq1$
\[
 F_{2n+1}^-\xrightarrow{\sim}F_{2n+2}^-\simeq F_3^{-\unsure{\varepsilon}}\otimes_k F_{2n-2}^+ \tag*{(21)}
\]
Choisissons
\[
 \delta\in t_E \tag*{(22)}
\]
et un
\[
 p\in F_2\qquad \varepsilon_2(p)=\delta^{\otimes2} \tag*{(23)}
\]
posons
\[
 p'=\Theta_\delta(p)\qquad\text{donc}\quad \varepsilon_3(p')=\delta^{\otimes3} \tag*{(24)}
\]
Alors
\[
 \left\{\begin{aligned}
 &1,p\quad\text{base de}\quad F_2=F_2^+=F_3^+\\
 &p'\quad\text{base de}\quad F_3^-\\
 &1,p,p'\quad\text{base de}\quad F_3=F_3^+\oplus F_3^-
 \end{aligned}\right. \tag*{(25)}
\]
\[
 \left\{\begin{aligned}
 &\{1,p,\ldots,p^n\}\ \text{base de}\ F_{2n}^+\simeq F_{2n+1}^+\\
 &\{p',p'p\ldots,p'p^{\unsure{n-1}}\}\ \text{base de}\ F_{2n+1}^-\simeq F_{2n+2}^-
 \end{aligned}\right. \tag*{$(25^1)$}
\]
\deleted{Propos.}
\[
 F_{2n+1}^-=F_{2n+2}^-=p'F_{2n-2}^+\qquad\text{pour }\deleted{\unsure{p|p'}}\,n\in\mathbb N^*. \tag*{(26)}
\]
\uline{Proposition} 1 Considérons l'homomorphisme

$f\longmapsto f\mid J$
\[
 F_3\longrightarrow k^J
\]
Alors on a
\[
 \operatorname{Ker}\bigl(F_3\mathrel{\unsure{-|\to}} k^J\bigr)=F_3^-=k.p' \tag*{(27)}
\]
% archive sheet number: 104
\scanpage{106}\mspage{3}
Si $f\in F_3^-$, plus gén. si $f\in F_n^-$, alors $x=-x$, pour $x\in E(k)$, $x\ne0$, implique $f(x)=-f(-x)=-f(x)$ i.e. $2f(x)=0$ i.e. $f(x)=0$. Ainsi $F_3\mathrel{\unsure{-|\to}}k^{J}$ est nul sur $F_3^-$, donc il reste a voir qu'il est injectif sur $F_3^+\simeq F_2$, Mais on a mieux

\uline{Corollaire} La matrice \deleted{formée} \inserted{de Vandermonde} $\bigl(p(i)^\alpha\bigr)_{\substack{i\in J\\0\leq\alpha\leq2}}$ est inversible, en d'autres termes les $p(i)$ $(i\in J)$ sont tous distincts.

OPS \deleted{\unsure{qu'on ait que pts de|qui sont les trois pts de}} \inserted{les pts de $J$} numérotés $0,1,\infty$, et $\Sigma_J=\mathbb P^1_k$. Quitte à multiplier par un scalaire non nul et ajouter une constante, ops
\[
 p(z)=\frac1{z-t}\qquad\text{où}\quad t=k\smallsetminus\{0,1\}=U_{0,3}(k)
\]
donc
\[
 p(0)=\frac1{-t},\qquad p(1)=\frac1{1-t},\qquad p(\infty)=0
\]
sont tous distincts, O.K.

\deleted{\unsure{$F$}} Considérons maintenant
\[
 \begin{aligned}
 &{p'}^{2}\in F_6^+=\deleted{\text{\unsure{somme directe de|somme directe des}}}\ \inserted{\text{espace ayant pour base}}\\
 &\hspace{5em}1,p,p^2,p^3
 \end{aligned}
\]
et notons que
\[
 \varepsilon_6{p'}^2=\varepsilon_6p^3=\delta^6
\]
on aura donc la relation
\[
 \boxed{{p'}^2=p^3+\deleted{\unsure{a|g}}a_1p^2+\deleted{\unsure{b}}a_2p+\deleted{\unsure{c}}a_3} \tag*{(28)}
\]
\margin{$a_1,\unsure{a_2},\unsure{a_3|s_3}\in k$ déterminés de façon unique (pour $p$ choisi)}
% archive sheet number: 105
\scanpage{107}
\uline{Proposition} 2 Les $a_\alpha$ ($1\leq\alpha\leq3$) sont caractérisés par la condition que $p^3+a_1p^2+a_2p+a_3\mid J=0$.

C'est nécessaire, puisque $p'\mid J=0$. Mais d'après le cor. à la prop. \unsure{précédent|précédente}, les fonctions $1,p,p^2$ restreintes à $J$ sont lin$^{\mathrm t}$ indépendantes, donc la relation précédente sur les $a_i$ (qui s'écrit
\[
 (-p^3)\mid J=a_1(p^2\mid J)+a_2(p\mid J)+a_3(1_E\mid J)\qquad)
\]
caractérise déjà les $a_i$, cqfd.

\uline{Corollaire 1} Les $\unsure{(-1)^\alpha/a_\alpha|\inserted{(-1)^\alpha}a_\alpha}$ ($1\leq\alpha\leq3$) sont les fonctions sym. élémentaires en les $p(i)$ ($i\in J$)

i.e. on a
\[
 p^3+a_1p^2+a_2p+a_3=\prod_{i\in J}\bigl(p-p\deleted{\unsure{f|\xi|'|p}}(i)\bigr)\,. \tag*{(29)}
\]
\deleted{\illegible{1}}
\uline{Corollaire 2} Pour que l'on ait $a_1=\unsure{0|g}$, il f. et s. que $\operatorname{Tr}_{J/k}(p\mid J)=0$ i.e. $\sum p(i)=0$.

Si car. $k\ne3$, les $p\in F_2$ qui satisfont cette condition forment un supplémentaire $W_E$ de $F_1$ dans $F_2$
\[
 F_2\simeq\underbrace{F_1}_{k.1_E}\oplus\underbrace{W_E}_{\simeq\, t_E^{\otimes2}} \tag*{(30)}
\]
Les \deleted{\unsure{sections|fonctions}} \inserted{éléments} de $W_E$ seront \inserted{donc} les fonctions de Weierstrass,
% archive sheet number: 106
\scanpage{108}\mspage{4}
en corr. $1-1$ avec les éléments de $(t_E^{\otimes2})^*$.

\uline{Corollaire 3} On a
\[
 \begin{aligned}
 F_\infty&\simeq k[P,P']\,/\,{P'}^2-(P^3+a_1P^2+a_2P^{\deleted{\unsure{2}}}+a_3)\\
 &\simeq k[P][P']\,/\,\bigl({P'}^2-(P^3+a_1P^2+a_2P^{\unsure{1|2}}+a_3)\bigr)
 \end{aligned} \tag*{(31)}
\]
l'isom étant donné par $P\mathrel{\unsure{-|\to}}p$, $P'\mathrel{\unsure{-|\to}}p'$. L'autom $f\to\check f$ correspond à l'autom. $P'\mapsto-P'$, $P\mapsto P$ de $k[P,P^{\unsure{+|'}}]/\ldots$ . La dérivation $\Theta_\delta$ est déduite par passage aux quotients de l'\deleted{application} dérivation
\[
 \Theta_{\deleted{\unsure{\delta}}}P=P',\qquad
 \Theta\deleted{\unsure{p}}P'=\frac12(3P^2+2a_1P+a_2) \tag*{(32)}
\]
sur $k[P,P']$, qui applique l'idéal $ {P'}^2-(P^3+a_1P^2+a_2P+a_3)$ en lui-$=$.

La formule (32) s'obtient en dérivant $ {p'}^2=p^3+a_1p^2+a_2p+a_3$ par $\Theta_\delta$, \deleted{et en divisant} \inserted{d'où}
\[
 2p'\Theta_\delta(p')=(3p^2+2a_1p+a_2)p'
\]
et en divisant par $p'$ la relation obtenue, d'où
\[
 p'=\tfrac12(3p^2+2a_1p+a_2)\,. \tag*{(33)}
\]
\deleted{\uline{Corollaire}}
\uline{Remarque} En termes des bases explicites plus haut des $F_n^+$ et $F_n^-$ $(25^1)$ donc des $F_n=F_n^+\oplus F_n^-$, on voit maintenant comment expliciter la loi de multiplication \inserted{en termes de $a_1,a_2,a_3$,} ainsi que
% archive sheet number: 107
\scanpage{109}
la dérivation $\Theta_\delta$.

Inversement, partant de $a_1,a_2,a_3\in k$, \deleted{satisfaisant} \inserted{tels que} le polynôme
\[
 Z^3+a_1Z^2+a_2Z+a_3\in k[Z]
\]
soit sans racines multiples, i.e tels que
\[
 \underbrace{\Delta(a_1,a_2,a_3)}_{\text{discriminant}}\ne0\,, \tag*{(34)}
\]
on vérifie que les formules précédentes définissent une courbe elliptique sur $k$ (qu'on peut \deleted{définir} \inserted{construire} comme le $\operatorname{Proj}(F_*)$, où $F_*$ est défini par (31)), \deleted{\unsure{dont voici}} \inserted{en tant} que rev. \deleted{\unsure{non ramifié}} \inserted{quadratique} de $\mathbb P^1_k$ défini par
\diagram{\begin{tikzcd}[ampersand replacement=\&,column sep=small,row sep=small]
 (35)\,\deleted{\unsure{p}}\qquad p:\ \& E \arrow[rr] \arrow[dr] \& \& \mathbb P^1_k\,\deleted{\unsure{\simeq}}\\
 \& \& E/\{\pm\operatorname{id}_E\}\arrow[ur,"\unsure{\sim}"] \&
\end{tikzcd}}
ramifié en \inserted{\unsure{l'ens.|l'im.} $J$} des trois points \inserted{$z_i$} de $\mathbb P^1_{(k)}\smallsetminus\{\infty\}$ \inserted{$=k$} solutions de l'équation
\[
 Z^3+a_1Z^2+a_2Z+a_3=0\,, \tag*{(36)}
\]
\deleted{la section} \inserted{l'élément} nul\deleted{\unsure{le|e}} de $E$ étant au dessus de \deleted{la section} \inserted{l'élément} $\infty$ de $\mathbb P^1_k$. \deleted{qui est} Bien \deleted{\unsure{que}} entendu, l'inclusion $J\subset\mathbb P^1_k$ définit un isom canonique
\[
 \Sigma_J\simeq\mathbb P^1_k,\qquad t\mapsto\infty\,. \tag*{(37}
\]
% archive sheet number: 108
\scanpage{110}\mspage{5}
La donnée d'un point $x$ \inserted{$\ne0$} de $E(k)$ équivaut à celle de son image dans $\mathbb P^1_k$, soit \deleted{$p(x)$}
\[
 p(x)=x_0\in k\,, \tag*{(38)}
\]
plus la donnée d'une des deux solutions \inserted{$x'$} de l'équation (\uline{deux} tout au moins si le $2^{\mathrm e}$ mb est $\ne0$ i.e. si $x\notin J={}_2E^*$ i.e. $x_0\notin D_J\subset\Sigma_J$)
\[
 {x'}^2=x_0^3+a_1x_0^2+a_2x_0+a_3\,, \tag*{(39)}
\]
donnée en termes de $x$ par
\[
 x'=p'(x)\,. \tag*{(40)}
\]
\uline{Proposition} 3 Soient $x,y,z\in E(k)\smallsetminus\{0\}$, \inserted{\deleted{pas} tous \deleted{égaux}\inserted{\uline{distincts}},} alors on a
\[
 x+y+z=0 \tag*{(\deleted{\unsure{4}}41)}
\]
SSS on a \deleted{a) On \unsure{n'a pas|a} $p(x)\unsure{=|\ne}p(y)=p(z)$ et b) on a}
\[
 \det\begin{pmatrix}
 1&p(x)&p'(x)\\
 1&p(y)&p'(y)\\
 1&p(z)&p'(z)
 \end{pmatrix}=0 \tag*{(42)}
\]
% slanted
\margin{\uline{NB} Si \deleted{deux} parmi les pts $x,y,z$ sont égaux, (42) est tjrs vérifié, sans qu'on ait \unsure{néc$^{\mathrm t}$} pour autant $x+y+z=0$ !}
% slanted
\margin{Cela implique qu'on a $p(x)\ne p(y)\ne p(z)\ne p(x)$ (car si on avait $p(x)=p(y)$ on aurait $x=y$ ou $x=-y$ d'où $z=0$), l'un et l'autre exclus}
% slanted
\margin{\deleted{La condition \unsure{a)|b)} \unsure{est vraie|peut même}, \unsure{car non a)|cas non a)} \unsure{signifierait|signifier} $\exists a\in E(k)$, avec deux \deleted{\illegible{1}} \deleted{\unsure{autres}} parmi les trois $x,y,z$ égaux à $a$, le troisième à $-a$, d'où $x+y+z=\unsure{a+a-a=a|a,\ \text{et}\ -a=a}$, donc $a=0$, ce qui est absurde puisque $x,y,z\ne0$.}}
En effet, la définition de la loi de multiplication de $E$ nous donne que si $x,y,z\in E(k)$ \unsure{autre condition|autre convention}, les \deleted{$x+y+z=0$}
\[
 \begin{aligned}
 x+y+z=0&\iff(\{x\}-\{0\})+(\{y\}-\{0\})+(\{z\}-\{0\})\overset{\rm lin.}{\sim}0\\
 &\text{i.e.}\quad\{x\}+\{y\}+\{z\}\overset{\rm lin.}{\sim}3\{0\}
 \end{aligned}
\]
Si $x,y,z$ sont tous $\ne0$, cela signifie qu'il existe une fonction rat. sur $E$, ayant $0$ comme pôle d'ordre 3, pas d'autre pôle, (donc $f\in F_3$), et $\{x\}+\{y\}+\{z\}$ comme \deleted{pôle} diviseur des zéros.
% archive sheet number: 109
\scanpage{111}
OPS $f\deleted{\unsure{p}}$ de la forme $p'+\alpha p+\beta$ (quitte à multiplier par un scalaire)
\[
 \left\{\begin{aligned}
 f&=p'+\alpha p+\beta,\quad \alpha,\beta\in k\\
 \operatorname{div} f&=\{x\}+\{y\}+\{z\}-3\{0_E\}
 \end{aligned}\right. \tag*{(43)}
\]
On aura donc
\[
 f(x)=f(y)=f(z)=0\qquad\text{i.e.} \tag*{(44)}
\]
\[
 \left\{\begin{aligned}
 p'(x)+\alpha p(x)+\beta&=0\\
 p'\unsure{y)|(y)}+\alpha p(y)+\beta&=0\\
 p'(z)+\alpha p(z)+\beta&=0
 \end{aligned}\right. \tag*{(45)}
\]
ce qui signifie aussi que le vecteur colonne $\begin{pmatrix}p'(x)\\p'(y)\\p'(z)\end{pmatrix}$ s'exprime comme comb. lin. des vecteurs colonnes $\begin{pmatrix}p(x)\\p(y)\\p(z)\end{pmatrix}$ et $\begin{pmatrix}1\\1\\1\end{pmatrix}$, \inserted{ce} qui \deleted{\unsure{don}}implique que ces trois vect. colonnes sont lin. dépendants, d'où (42). Inversement, supposons qu'on ait (42), et $x,y,z$ \deleted{\unsure{deux|égaux}} \inserted{tous distincts}, \deleted{\unsure{par ex|pas}} \deleted{que l'on} \inserted{alors on} n'a pas $p(x)=p(y)=p(z)$, \deleted{\unsure{sinon}} \inserted{(sinon $x,y,z$ appartiendraient à une $=$ fibre de $p$, donc deux parmi eux seraient égaux)}, \deleted{\unsure{et que}} \deleted{\unsure{alors}} \inserted{donc} les vect. colonne $\begin{pmatrix}1\\1\\1\end{pmatrix}$ $\begin{pmatrix}p(x)\\p(y)\\p(z)\end{pmatrix}$ sont lin. indép., et (42) signifie qu'il existe $\alpha,\beta\in k$ tels qu'on ait (45), i.e. (44), où $f$ est donnée par (43) a). Comme $x,y,z$ sont distincts, on en conclut bien (43) b), d'où (41).
% archive sheet number: 110
\scanpage{112}\mspage{6}
Montrons comment cela permet de calculer $z$, via $z_0=p(z)$ et $z'=p'(z)$ (reliés par
\[
 {z'}^2=z_0^3+a_1z_0^2+a_2z_0+a_3\qquad), \tag*{(46)}
\]
quand on connaît $x$ et $y$, exprimés à l'aide de \deleted{$x_0=p(x),x'=p'(x)$}
\[
 \left\{\begin{aligned}
 x_0&=p(x),\ x'=p'(x)\\
 y_0&=p(y),\ y'\mathrel{\unsure{-|=}}p'\unsure{y)|(y)}
 \end{aligned}\right.
\]
(reliés par
\[
 {x'}^2=x_0^3+a_1x_0^2+a_2x_0+a_3,\quad
 {y'}^2=y_0^3+a_1y_0^2+a_2y_0+a_3).
\]
L'équation (\unsure{41|4)|42}) s'écrit
\[
 z'(x_0-y_0)+z_0(y'-x')+(x'y_0-y'x_0)=0 \tag*{$(*)$}
\]
i.e.
\[
 z'=\lambda z_0+\mu,\quad\text{avec}\quad
 \lambda=\frac{x'-y'}{x_0-y_0},\quad
 \mu=\frac{y'x'_0-x'y_0}{x_0-y_0} \tag*{(47)}
\]
\deleted{Donc} (NB on a $x_0\ne y_0$ i.e $x_0-y_0\ne0$).

De plus $z'$ satisfait (46), qui équivaut \deleted{à} \deleted{\unsure{l'équa|l'équation}} (compte tenu de (47), qui implique $ {z'}^2=\lambda^2z_0^2+2\lambda\mu z_0+\mu^2$) à l'équation
\[
 z_0^3+(a_1-\lambda^2)z_0^2+(a_2-2\lambda\mu)z_0+(a_3-\mu^2)=0. \tag*{(48)}
\]
Donc $(z_0,z')$ est lié par (47), (48), a priori on trouve 3 solutions \inserted{en $z_0$} de (48), à qui correspondent trois solutions $(z_0,z')$ via (47). Il faut trouver une solution unique, on se rappelle qu'on doit avoir
\[
 z_0\ne x_0,\quad z_0\ne y_0,\qquad\deleted{\text{\unsure{do|de}}} \tag*{(49)}
\]
% archive sheet number: 111
\scanpage{113}
alors qu'il est clair que $\inserted{\unsure{z_0}=}x_0$, $\inserted{z_0=}y_0$ sont solutions solutions de l'équation (48). Si donc $z_0$ est une solution $\ne x_0,y_0$, on aura
\[
 \begin{aligned}
 x_0+y_0+z_0&=-(a_1-\lambda^2)=\lambda^2-a_1\\
 &=\left(\frac{x'-y'}{x_0-y_0}\right)^2\unsure{-|=}a_1
 \end{aligned}
\]
d'où la solution (sous réserve d'existence d'une solution $z_0\ne x_0,y_0$, ou encore \deleted{\unsure{d'une so|d'autre part}} \inserted{sous} réserve que
\[
 \left\{\begin{aligned}
 &x,y\in E\smallsetminus\{0\},\quad x\ne y,\quad -(x+y)\ne x,y\quad\text{i.e.}\\
 &\hspace{10em}2x+y\ne0,\quad 2y+x\ne0\\
 &\hspace{10em}\text{i.e.}\quad y\ne-2x,\ x\ne-2y\quad),
 \end{aligned}\right. \tag*{(50)}
\]
la solution
\[
 \left\{\begin{aligned}
 z_0&=-x_0-y_0+\left(\frac{x'-y'}{x_0-y_0}\right)^2-a_1\\
 z'&=\lambda z_0+\mu=\frac{x'-y'}{x_0-y_0}
 \left[-x_0-y_0+\left(\frac{x'-y'}{x_0-y_0}\right)^2-a_1\right]\\
 &\hspace{9em}\inserted{+\frac{y'x_0-x'y_0}{x_0-y_0}}\,.
 \end{aligned}\right. \tag*{(51)}
\]
Par principe de prolongement, cette expression pour $z_0=p(z)$ et $z'=p'(z)$, pour $z=-(x\unsure{-|+}y)$, est valable encore dès que $x,y\in E(k)\smallsetminus\{0\}$ sont tels que $x_0\ne y_0$ i.e. $p(x)\ne p(y)$, i.e. dès que $\deleted{x\unsure{=|\ne}}y\notin\pm x$, de sorte que les formules (51) ont un sens.
% archive sheet number: 112
\scanpage{114}
La fonction $p$, considérée comme morphisme
\[
p:E\longrightarrow\mathbb{P}^1_k\tag*{(52)}
\]
précise l'iso
\[
\dot p:E/\{\pm\mathrm{id}_E\}\simeq\text{droite projective type }\mathbb{P}^1_k\tag*{(53)}
\]
par la condition
\[
p(0_E)=\infty\tag*{(54)}
\]
ce qui détermine $p$ modulo transformation affine
\[
p\longmapsto p_1=ap+b\qquad a\in k^*,\ b\in k.\tag*{(55)}
\]
Les points d'ordre $2$
\[
(\unsure{\varepsilon|\xi}_i)_{i\in J},\quad \varepsilon_i\in{}_{2}E(k)^*\tag*{(56)}
\]
sont transformés par $p$ en des points
\[
\zeta_i\in p(\varepsilon_i)\quad\in\mathbb{P}^1_{\deleted{\unsure{k|*}}}(k)\setminus\{\unsure{\infty|0}\}=k\tag*{(57)}
\]
% slanted
\margin{\uline{NB} On obtient, via des $E,p$, \unsure{n'importe} quel triplet de points distincts de $k$}
Ce sont, bien sûr, avec le point $\infty$, les points de ramification de $p$. La normalisation de Weierstrass, possible en car. $\ne3$, \deleted{qui} est
\[
\underbrace{\sum\zeta_i}_{-a_1}=0\tag*{(58)}
\]
elle détermine $p$ modulo transformation
\[
p\longmapsto p_1=ap\qquad a\in k^*.
\]
En général, on pose
\[
\left\{\begin{aligned}
a_1&=-\sum\zeta_i\\
a_2&=\sum\zeta_i\zeta_{\unsure{1|j}}\\
a_3&=\deleted{\unsure{\prod\zeta_i|\zeta_i}}-\prod\zeta_i
\end{aligned}\right.\tag*{(59)}
\]
% archive sheet number: 113
\scanpage{115}
de sorte que $\zeta_1,\zeta_2,\zeta_3$ sont les solutions de l'équation
\[
\underbrace{z^3+a_1z^2+a_2z+a_3}_{\substack{\deleted{\unsure{(z-\zeta_1)(z-\zeta_2)|(z-\zeta_j)(z-\zeta_k)}}\\\prod_i(z-\zeta_i)}}=0\tag*{(60)}
\]
On aura, \deleted{\unsure{donc|dans}} si
\[
\delta\in t_E=T_{E,0}\quad\deleted{\unsure{\delta|\xi}}
\]
est une base de $t_E$, liée à $p\bmod k.1$ par
\[
\varepsilon_2(p)\simeq\delta^{\otimes2}
\]
et si
\[
p'=\Theta_\delta p
\]
la relation
\[
p'{}^2=p^3+a_1p^{\unsure{2|7}}+a_2p+a_3.
\]

Déterminons l'autom. involutif $\sigma_i$ de $\mathbb{P}^1_k$ qui échange $\zeta_i$ et $\infty$, $\zeta_j$ et $\zeta_k$ -- qui \deleted{\unsure{est donc|est donné}}\inserted{s'identifie} à l'autom. déduit par passage au quotient de
\[
\deleted{x\longmapsto x+}\qquad x\longmapsto x+\varepsilon_i\qquad E\longrightarrow E
\]
On trouve, \deleted{\unsure{utilisant|utilisons}}\inserted{écrivant,}
\[
\sigma_i z=\frac{az\unsure{-|+}b}{cz\unsure{-|+}d}
\]
et écrivant $\sigma_i\infty=\zeta_i,\ \sigma_i\zeta_i=\infty$
\[
\deleted{\sigma_i z}\quad\frac ac=\zeta_i,\quad-\frac dc=\zeta_i,
\]
soit, faisant $c=1$
\[
\sigma_i z=\frac{\zeta_i z+b}{\deleted{\unsure{\zeta_i|a_i}}z-\unsure{\zeta_i|\xi_i}}
\]
et écrivant $\sigma_i\zeta_j=\zeta_k$ on trouve
\[
\frac{\unsure{\zeta_i|\xi_i}\zeta_j+b}{\zeta_j-\zeta_i}=\zeta_k\qquad\text{i.e.}\qquad b=\zeta_j\zeta_k-\zeta_i(\zeta_j+\zeta_k)
\]
% archive sheet number: 114
